\documentclass[10pt,a4paper]{amsart}
\usepackage[margin=19mm]{geometry}
\usepackage{amsmath,amssymb,mathtools}
\usepackage[scr=rsfs,cal=euler]{mathalfa}
\usepackage{xfrac}
\usepackage[unicode,hidelinks,bookmarks=false]{hyperref}
\newcommand{\ie}{i.e.\ }
\newcommand{\cf}{cf.\ }
\newcommand{\resp}{respectively\ }
\newcommand{\Subsection}{\subsection}
\providecommand{\nobreakdash}{\nobreak\hbox{-}}
\setlength{\emergencystretch}{2em}
\multlinegap0pt
\usepackage{bm}
\usepackage{bbm}
\usepackage{dsfont}
\usepackage{xspace}
\usepackage{cancel}
\usepackage{mathtools}
\usepackage{cleveref}
\usepackage{amsthm}
\makeatletter
\@ifundefined{dummy}{\newtheorem{dummy}{Dummy}}{}
\@ifundefined{lemma}{\newtheorem{lemma}[dummy]{Lemma}}{}
\@ifundefined{theorem}{\newtheorem{theorem}[dummy]{Theorem}}{}
\makeatother
\DeclareMathOperator{\Lie}{Lie}
\DeclareMathOperator{\spn}{span}
\newcommand{\E}{\mathcal{E}}
\newcommand{\LM}{\mathcal{M}}
\newcommand{\R}{\mathbb{R}}
\newcommand{\s}[1]{\Gamma(#1)}
\title[Controllability on landmark manifolds for shapes and neural ]{Controllability on landmark manifolds for shapes and neural ODEs}
\author{Erlend Grong, Sylvie Vega-Molino}
\date{Source: arXiv:2403.08090v3 (CC BY 4.0)}
\begin{document}
\maketitle

\noindent\emph{Source and licence.} Controllability on landmark manifolds for shapes and neural ODEs. Version arXiv:2403.08090v3 (CC BY 4.0), distributed under the Creative Commons Attribution 4.0 International licence, as recorded in the arXiv licence metadata for this exact version and shown on the abstract page https://arxiv.org/abs/2403.08090v3. Full rights notice: licenses/arxiv-2403.08090-CC-BY-4.0.txt.\par\smallskip

\noindent\textit{Editorial scope.} Counted unit: the Theorem whose source label is \texttt{th:Higher} in the article (source0.tex lines 407--414, source label \texttt{th:Higher}), delivered together with the article's own complete original proof of it (source0.tex lines 473--550, one \texttt{proof} environment of 78 lines). No mathematics is written, repaired or re-derived: statement and proof are the article's own text line for line. The proof is detached from the statement in the source: the article states the result at lines 407--414 and prints its proof at lines 473--550, and the proof environment's own header names the statement, so the association is the article's own. Required context retained uncounted: lemma:polynomials (lines 437--441). Every definition, display and label of those lines is retained unchanged. Omitted from this package and not redistributed: 1--406, 415--436, 442--472, 551--562. Nothing inside a retained excerpt is omitted or altered: every retained body is delivered exactly as the article's lines are written, so no omission receipt of any kind accompanies this package. Citations: the retained excerpts carry no citation command at all, so no bibliography entry of the article is transcribed and no reference is redistributed. Reference adaptation: none was needed and none was made; every reference inside the delivered unit resolves inside it, so no unresolved reference or citation is produced. Printed slips kept literal: no printed word, symbol, hypothesis or number of the counted statement or of its proof was altered. Layout and class: the article's own class is \texttt{amsart} with packages including amsmath, amsfonts, amssymb, amsthm, amsbsy, xcolor, mathrsfs, enumitem, todonotes, soul, hyperref, cleveref, autonum; the wrapper is the lane's pinned \texttt{amsart} editorial wrapper at 10pt on A4, which restates the theorem-like environments the retained text needs and adds the article's own macro definitions that the retained text needs (\texttt{\textbackslash E}, \texttt{\textbackslash LM}, \texttt{\textbackslash Lie}, \texttt{\textbackslash R}, \texttt{\textbackslash s}, \texttt{\textbackslash spn}), with extra wrapper packages (bm, bbm, dsfont, xspace, cancel, mathtools, cleveref). The delivered numbering is the unit's own. Assets: no retained span contains a figure environment, a graphics-inclusion command, a PDF-page inclusion, a table or a TikZ picture; no image file is copied into this package, the package declares no asset dependency and no image file is redistributed with it. Licence: version arXiv:2403.08090v3 of this article is distributed under the Creative Commons Attribution 4.0 International licence, as recorded in the arXiv licence metadata for this exact version and shown on the abstract page https://arxiv.org/abs/2403.08090v3. A full rights notice is delivered at licenses/arxiv-2403.08090-CC-BY-4.0.txt. Characters: every retained excerpt is pure ASCII, so the delivered engine, XeLaTeX with its default Latin Modern text font, reports no missing character.
\begin{lemma} \label{lemma:polynomials} Consider the landmark manifold $\LM = \LM^n(\mathbb{R}^d)$.
Let $\E$ be a subbundle of $T\LM$ with $\Lie(\E)$ denoting all vector fields that can be generated from iterated brackets of its sections. Assume that for any $j=1,\dots, d$ and for any polynomial $p$ on $\mathbb{R}^d$ of degree equal to $n-1$,
$$\ell (p \partial_j) \in \Lie(\E).$$
Then $\E$ is bracket-generating.
\end{lemma}

\begin{theorem} \label{th:Higher}
Consider the vector field on $\mathbb{R}^d$ for $d \geq 3$,
\begin{equation}
Y = \sum_{k=1}^d \left( \sum_{j=1}^d (x^j)^2 \right) x^k \partial_{k+1} =\| x\|^2 \sum_{k=1}^d  x^k \partial_{k+1} \in \s{T\R^d}
\end{equation}
If $\E= \spn \{ \ell \partial_1, \ell Y\} = \E \subseteq T\LM$,
then we have complete controllability with respect to $\E$ independent of $n$.
\end{theorem}

\begin{proof}[Proof of \Cref{th:Higher}]
We will first show that that we can use Lie brackets to generate any vector field on the form $(x^j)^\alpha \partial_k$ from $\partial_1$ and $Y$ for $j,k =1, \dots, d$, $\alpha = 0,1, 2, \dots$. We will then show that this is sufficient to generate all polynomials.

Recall that indices are denoted modulo $d$. Observe that
\begin{align}
[\partial_j, Y] &= 2 x^j \sum_{k=1}^d x^k \partial_{k+1} + \|x\|^2 \partial_{j+1}, \\
[\partial_j, [\partial_j, Y]] &= 2 \sum_{k=1}^d x^k \partial_{k+1} + 4 x^j \partial_{j+1}  ,\\
[\partial_j, [\partial_j, [\partial_j, Y]]] &= 6\partial_{j+1}.
\end{align}
Therefore, starting with $\partial_1$, it follows that we can inductively generate all degree 0 monomials vector field $\partial_j$.

For degree one vector fields, observe that for $j \neq k$,
\begin{align}
\frac{1}{2} [\partial_j, [\partial_k, Y]] &=  x^k  \partial_{j+1} + x^j \partial_{k+1} ,\\
\frac{1}{4}([\partial_j, [\partial_j, Y]]-[\partial_k, [\partial_{k}, Y]]]) &=  x^j \partial_{j+1} - x^k \partial_{k+1}.
\end{align}
Since we are assuming $d >2$, we have $j+2 \neq j \bmod d$, and hence we can use
$$[x^{j+1} \partial_{j+1} + x^j \partial_{j+2}, x^{j+2} \partial_{j+1} + x^{j} \partial_{j+3}]  = (x^j  - x^{j+2}) \partial_{j+1}.$$
Furthermore, we observe that
$$[x^{j-1} \partial_{j+1} + x^j \partial_j, (x^j - x^{j+2}) \partial_{j+1}] = x^j \partial_{j+1},$$
giving us all monomials of the form $x^j\partial_{j+1}$. Combining these vector fields, we first obtain $[x^{j} \partial_{j+1}, x^{j+1} \partial_{j+2}] = x^j \partial_{j+2}$ and iterating this bracket, we are able to generate vector fields of the form $x^l \partial_{l}$, $j \neq l$ with our Lie brackets. Finally,
\begin{equation}
[\partial_{j-1}, [\partial_j, Y]] = 2(x^{j-1}\partial_{j+1} + x^j\partial_j)
\end{equation}
and so by linear combination with previously acquired terms, we obtain vector fields, $x^j \partial_{j}$, giving us a complete set of all degree one monomials $x^j\partial_l$, $j,l=1,\dots, d$.

In degree two, we consider the brackets
\begin{align}
[x^j \partial_j, [\partial_j, Y]] &= 2 x^j \sum_{k=1}^d x^k \partial_{k+1} + 4 (x^j)^2 \partial_{j+1} - 2 x^j x^{j-1} \partial_{j} \\
& = 6(x^j)^2 \partial_{j+1} + \sum_{l \neq j, j-1} 2x^j x^l \partial_{l+1} , \\
[x^j \partial_j, [x^j \partial_j, [\partial_j, Y]]] & = 12(x^j)^2 \partial_{j+1} + \sum_{l \neq j, j-1} 2x^j x^l \partial_{l+1} ,
\end{align}
from which we can take the difference to generate all terms of the form $(x^j)^2 \partial_{j+1}$.  For $j \neq k+1$ it holds
\begin{align} \label{Special}
[x^j\partial_k, (x^k)^2\partial_{k+1}] &= 2x^kx^j \partial_{k+1} ,\\
[x^j\partial_k, [x^j\partial_k, (x^k)^2\partial_{k+1}]] &= 2(x^j)^2 \partial_{k+1} ,
\end{align}
giving us all monomials $(x^j)^2 \partial_k$ with $k \neq j$. To get the final second order monomial, we note that
\begin{align}[x^k \partial_j, (x^j)^2 \partial_k] &= 2x^k x^j \partial_k - (x^j)^2 \partial_j, \\
[x^k \partial_{j}, x^{k-1} x^j \partial_k] &= x^k x^{k-1} \partial_k - x^{k-1} x^j \partial_{j} ,
\end{align}
and combining these identities with \eqref{Special}, and using linear combinations, we are finally able to generate monomials $(x^j)^2 \partial_j$, giving us all degree two monomials $(x^j)^2\partial_k$ for $j,k = 1, \dots, d$.

In degree three, it follows that for $j,k,l$ all different
\begin{align}
[(x^k)^2\partial_l, (x^l)^2 \partial_j] &= 2(x^k)^2 x^l \partial_j ,\\
[x^k \partial_l, [(x^k)^2\partial_l, (x^l)^2 \partial_j] &= 2(x^k)^3 \partial_j,
\end{align}
For identical coefficients, we first compute
% \begin{align}
% & [x^{j+1} \partial_j ,[x^{j+1} \partial_j,[x^{j+1} \partial_j, \sum (x^l)^2 x^k \partial_{k+1}]]] \\
% & = 2[x^{j+1} \partial_j ,[x^{j+1} \partial_j, \sum x^j x^{j+1} x^k \partial_{k+1}]]\\
% & \qquad + [x^{j+1} \partial_j,[x^{j+1} \partial_j, \sum (x^l)^2 x^{j+1} \partial_{j+1}]] \\
% & \qquad - [x^{j+1} \partial_j,[x^{j+1} \partial_j, \sum (x^l)^2 x^j \partial_{j}]] \\
% & = 2 [x^{j+1} \partial_j , \sum (x^{j+1})^2 x^k \partial_{k+1}]+ 4 [x^{j+1} \partial_j ,  x^j (x^{j+1})^2 \partial_{j+1}]\\
% & \qquad - 4[x^{j+1} \partial_j, (x^j)^2 x^{j+1} \partial_{j}]  - [x^{j+1} \partial_j, \sum (x^l)^2 x^{j+1} \partial_{j}] \\
% & = 6 (x^{j+1})^3 \partial_{j+1} - 16x^j (x^{j+1})^2 \partial_{j}   \\
% \end{align}
\begin{align}
[x^{j+1}\partial_j, [x^{j+1}\partial_j, [x^{j+1}\partial_j, Y]]] &= 6(x^{j+1})^3 \partial_{j+1} - 16(x^{j+1})^2 x^j \partial_j,
\end{align}
and then
\begin{equation}
[x^{j+1}\partial_j, [x^{j+1}\partial_j, [x^{j+1}\partial_j, Y]]] + 8 [(x^{j+1})^2\partial_j, (x^j)^2 \partial_j] = 6(x^{j+1})^3 \partial_{j+1}.
\end{equation}
In total, we obtain all pure degree three monomials $(x^j)^3 \partial_k$.

For $\alpha \geq 3$ it holds that
\begin{equation}
[(x^i)^2\partial_i, (x^i)^\alpha \partial_j] = \begin{cases}
\alpha(x^i)^{\alpha+1} \partial_j & i \neq j \\
(\alpha-2)(x^i)^{\alpha+1} \partial_i & i = j \\
\end{cases}
\end{equation}
from which we generate all pure monomials $(x^j)^\alpha \partial_i$. Finally, we show that we can obtain arbitrary polynomials. Any polynomial vector field can be written as a sum of monomials
$$(x^{i_1})^{\alpha_1} \cdots (x^{l})^{\alpha_l} (x^j)^{N - l} \partial_j, \qquad i_r \neq j, r=1, \dots, l$$
and we can obtain such a vector field by starting with $(x^j)^N \partial_j$ and then in order take the bracket with vector field $(x^{i_r})^{\alpha_r} \partial_j$ for $j=1, \dots l$. The result now follows from \Cref{lemma:polynomials}.
\end{proof}

\end{document}
